\cotefile{84}
\scanpage{51}
% manuscript line(s) 1
\[
N=(T^2\deleted{\unsure{+}}\unsure{-}\delta)/4,
\]
% manuscript line(s) 2–3
c'est une section de $\Gamma\,\mathrm{Sym}^2(E)$, grâce à la condition $T_0^2\equiv\delta\ (4)$. Vérifions que \unsure{le|la} \deleted{fonction} pol.
% manuscript line(s) 4
\[
Q(x)=x^2-Tx+N
\]
% manuscript line(s) 5–6
sur $P_1$ est à valeurs \deleted{sur} $L^{\otimes2}\subset\mathrm{Sym}^2(E)$. Il suffit de prouver qu'il en est ainsi de
% manuscript line(s) 7
\[
4Q(x)=Q(2x)
\]
% manuscript line(s) 8
et on peut poser
% manuscript line(s) 9
\[
2x=T+z
\]
% manuscript line(s) 10
on trouve
% manuscript line(s) 11
\[
Q(T+z)=T^2+2zT+z^2-\underset{\substack{\shortparallel\\\psi(T+z)}}{2}T(T+z)+\underset{\substack{\shortparallel\\\psi(T+z)^2}}{4}N\qquad\text{i.e.}
\]
% manuscript line(s) 12
\[
\boxed{Q(T+z)=\deleted{\unsure{4T^2+}}\,z^2\deleted{\unsure{+}}\unsure{-}\delta}
\]
% manuscript line(s) 13
et on a bien $z^2\deleted{\unsure{+}}\unsure{-}\delta\in\Gamma L^{\otimes2}$ \quad OK.

% manuscript line(s) 14
Description de $\mathcal A$. \unsure{On|Or|Car|On a},
% manuscript line(s) 15
\[
\mathcal A\longrightarrow(\underline O\times L^\vee)\qquad x\longrightarrow(\operatorname{Tr}(x),\varphi(x))
\]
% manuscript line(s) 16–17
est injectif, car le noyau est formé des $x\in\Gamma(\underset{\substack{\cap\\\mathcal A}}{\underline O})$ tels que $\operatorname{Tr}x=0$ i.e. $2x=0$ i.e. $x=0$.
% manuscript line(s) 18
Donc
% manuscript line(s) 18
\[
\mathcal A\subset\underline O\times L^\vee
\]
% manuscript line(s) 19
On trouve
% manuscript line(s) 19
\[
\Gamma\mathcal A\simeq\left\{(b,\unsure{r|\gamma})\in\Gamma\underline O\times\Gamma L^\vee\ \left|\ \boxed{\deleted{\unsure{r_0}}b_0=\deleted{\unsure{b_0}}r_0T_0}\right.\right\}
\]
% manuscript line(s) 20
i.e. $b\in\Gamma\underline O_X$ relève $\underset{\substack{\in\\\Gamma L_0^\vee}}{\deleted{\unsure{b_0}}r_0}\cdot\underset{\substack{\in\\\Gamma L_0}}{T_0}\in\Gamma\underline O_{X_0}$.
% manuscript line(s) 21–22
\diagram{\begin{tikzcd}[ampersand replacement=\&,column sep=large] \unsure{0}\arrow[r]\&{\underline O_X}^{\unsure{0}}\arrow[r,"\xi\mapsto2\xi"]\&\mathcal A\arrow[r,"{(b,r)\mapsto r}","\mathrm{surj}"']\&L^\vee\arrow[r,dash]\&0\end{tikzcd}}
% manuscript line(s) 23
Ainsi l'élém\textsuperscript{t} $\uline{1}$ de $\mathcal A$ est décrit par $(2,0)$
% archive sheet number: 50
\scanpage{52}
% manuscript line(s) 1–3
\[
\deleted{\unsure{N(b|\mathrm{N.B.}}}\quad\begin{aligned}&\left\{\begin{aligned}T(b,\unsure{r|\gamma})&=b\\ \delta(b,r)&=\deleted{\unsure{b}}r^2\cdot\delta\\ N(b,r)&=(T^2\deleted{\unsure{+}}\unsure{-|+}\delta)(b,r)/4=(b^2\deleted{\unsure{+}}-r^2\cdot\delta)/4\end{aligned}\right. \end{aligned}
\]
% manuscript line(s) 3
\margin{donc}
% manuscript line(s) 4
C'est une section de \deleted{\unsure{$\Gamma$|de}}$\underline O_X$, car
% manuscript line(s) 5
\[
b\equiv rT\ (2)\qquad\text{donc}\qquad b^2\equiv r^2T^2\ (4)
\]
% manuscript line(s) 5–7
% slanted
% manuscript line(s) 5–7
\margin{\uline{NB} On en déduit $\varphi_N((b,r),(b',r'))=\dfrac{bb'\deleted{\unsure{+}}\inserted{-}rr'\cdot\delta}{\deleted{\unsure{4}}2}$}
% manuscript line(s) 6
donc
% manuscript line(s) 7
\[
b^2\deleted{\unsure{+}}-r^2\cdot\delta\unsure{\simeq|\equiv}r^2(\underbrace{T^2\deleted{\unsure{+}}-\delta}_{\equiv0\ (4)})\equiv0\ (4)
\]

% manuscript line(s) 8
Soit $\underline{x\in(b,r)}\in\Gamma\mathcal A$, on a
% manuscript line(s) 9
\[
x^2\deleted{\unsure{+}}-bx\deleted{\unsure{-}}+\frac{b^2\deleted{\unsure{+}}-r^2\cdot\delta}{4}=0\qquad\text{i.e.}
\]
% manuscript line(s) 10–11
\[
\begin{aligned}x^2&=bx-\unsure{\cdot}\frac{b^2-r^2\cdot\delta}{4}=\left(b^2-\frac{b^2-r^2\delta}{2},br\right)\\ &\phantom{=bx-\frac{b^2-r^2\cdot\delta}{4}}=\left(\frac{b^2+r^2\cdot\delta}{2},br\right)\end{aligned}
\]
% manuscript line(s) 12
Cela permet de calculer, si
% manuscript line(s) 13
\[
x=(b,r),\qquad x'=(b',r')
\]
% manuscript line(s) 14
le produit $xx'$ comme
% manuscript line(s) 15
\[
xx'=\tfrac12\big((x+x')^2-x^2-x'^2\big)
\]
% manuscript line(s) 16–18
\[
\begin{aligned}\text{i.e.}\quad 2xx'=\Big(&\frac{(b+b')^2+(r\unsure{-|+}r')^2\cdot\delta}{2}-\deleted{\unsure{b^2}}\frac{b^2+r^2\cdot\delta}{2}-\frac{b'^2+r'^2\cdot\delta}{2},\\ &(b+b')(r+\deleted{\unsure{b'}}r')\deleted{\unsure{r}}-br-b'r'\Big)\Big)\\ =\Big(&\underbrace{bb'+rr'\cdot\delta}_{\substack{\mapsto r_0r'_0\delta_0+r_0r'_0\delta_0\\=2r_0r'_0\delta_0=0}},\underset{\substack{r_0r'_0T_0+r_0r'_0T_0\\=2r_0r'_0T_0=0}}{br'+b'r}\Big)\end{aligned}
\]
% manuscript line(s) 18–20
\margin{mod $2$ c'est bien nul car $b_0=r_0T_0$, $b'_0=r'_0T_0$ $T_0^2=\delta_0$}
% manuscript line(s) 20
donc
% manuscript line(s) 21
\[
\boxed{xx'=(b,r)(b',r')=\left(\tfrac12(bb'\deleted{\unsure{-}}+rr'\cdot\delta,\tfrac12(br'+rb')\right)\big)}
\]
% archive sheet number: 51
\scanpage{53}

% manuscript line(s) 1–2
L'accouplement $\mathcal A\times E\longrightarrow\underline O$ \quad $(x,f)\longmapsto\langle x,f\rangle=f(x)$ est donné par
% manuscript line(s) 3
\[
\boxed{\langle(b,\unsure{r|\gamma}),(\lambda,\unsure{\zeta|\xi})\rangle=\tfrac12(\lambda b+r\cdot\zeta)}
\]
% manuscript line(s) 4
C'est OK car
% manuscript line(s) 5
\[
\underbrace{\lambda_0\underset{\substack{\shortparallel\\r_0\cdot T_0}}{b_0}}_{\lambda_0r_0\cdot T_0}+\underbrace{r_0\cdot\underset{\substack{\deleted{\unsure{r_0T_0}}\shortparallel\\\lambda_0T_0}}{\zeta_0}}_{r_0\lambda_0T_0}=2r_0\lambda_0T_0=0
\]

% manuscript line(s) 6
On a
% manuscript line(s) 7–8
\[
\begin{aligned}\varphi_N\big((\underset{\substack{\in\\\Gamma O}}b,\underset{\substack{\in\\\Gamma L^\vee}}r),(\underset{\substack{\in\\\Gamma\underline O}}{b'},\underset{\substack{\in\\\Gamma L^\vee}}{r'},)\big)&=\tfrac12(bb'\deleted{\unsure{+}}\unsure{-|+}rr'\cdot\delta)\\ &=\langle(\underset{\substack{\in\\\Gamma O}}b,\underset{\substack{\in\\\Gamma L^\vee}}r),(\underset{\substack{\in\\\Gamma O}}{b'},\underbrace{\unsure{-}r'\delta}_{\in\Gamma L})\rangle\end{aligned}
\]
% manuscript line(s) 8
, donc
% manuscript line(s) 9
$f_N:\mathcal A\unsure{\mathbin{-}|\to}E$ est donné par
% manuscript line(s) 10
\[
\boxed{f_N(\underset{\substack{\in\\\Gamma\underline O}}b,\underset{\substack{\in\\\Gamma L^\vee}}r)=(\underset{\substack{\in\\\Gamma\underline O}}b,\underset{\substack{\in\\\Gamma L}}{\deleted{\unsure{+}}\inserted{-}r\delta})}
\]

% manuscript line(s) 11–12
\deleted{Vérifions} \inserted{Utilisons} la formule $Q(f_N(x))=+N(x)\delta$ par la formule précédente, pour $f_N(x)$, $x=(b,r)$, \deleted{\unsure{être}}.
% manuscript line(s) 13–15
\deleted{\unsure{Q}} donc
% manuscript line(s) 13–15
\[
\begin{aligned}Q(\underset{\substack{\shortparallel\\\lambda}}{b},\underbrace{-r\cdot\delta}_{\shortparallel\,\zeta})&=\deleted{\unsure{-}}(b^2-r^2\cdot\delta)/4\ \delta\\ &=\deleted{\unsure{-}}(\lambda^2+r\zeta)/4\ \delta\\ &\text{où}\qquad r=\unsure{-}\zeta\cdot\delta^{-1}\end{aligned}
\]
% manuscript line(s) 16
d'où
% manuscript line(s) 16–17
\[
\begin{aligned}Q(\lambda,\zeta)&=+(\lambda^2\deleted{\unsure{-}}\unsure{+|-}\zeta^2\delta^{-1})/4\ \delta\\ &=\deleted{\unsure{\tfrac14}}(\zeta^2-\lambda^2\delta)/4\end{aligned}
\]
% manuscript line(s) 18
i.e.
% manuscript line(s) 18
\[
\boxed{Q(\lambda,\zeta)=(\zeta^2\deleted{\unsure{+}}-\lambda^2\delta)/4}
\]
% manuscript line(s) 19
Il faut vérifier que $\zeta^2-\lambda^2\delta\equiv0\ (4)$
% manuscript line(s) 19–21
\margin{Or $\zeta\equiv\lambda T(2)$ d'où \deleted{\unsure{la}} $\zeta^2\equiv\lambda^2T^2\ (4)$ donc $\zeta^2-\lambda^2\delta\underset{\unsure{4|d}}\equiv\lambda^2(\underbrace{T^2-\delta}_{4N})$}
% manuscript line(s) 20–21
$=4\lambda^2N\equiv0\ (4)$ OK \deleted{$\zeta_0^2=\lambda_0^2\delta_0$ Or on a $\zeta_0=\lambda_0T_0$ $\zeta_0^2=\lambda_0^2T_0^2=\lambda_0^2\delta_0$ car $T_0^2=\delta_0$ \unsure{OK!}}
% archive sheet number: 52
\scanpage{54}

% manuscript line(s) 1–2
Si $y=(\lambda,\unsure{\zeta|\xi})\in\Gamma P_1$ \deleted{\unsure{alors}} \inserted{i.e. $\lambda=1$,} l'équation $Q(y)=0$ équivaut bien à $(\zeta^2\deleted{\unsure{+}}\unsure{-|+}\delta)/4=0$ i.e.
% manuscript line(s) 3
\[
\zeta^2=\delta
\]
% manuscript line(s) 4–7
(en plus on a $\zeta_0=T_0$, sous-entendu), cela prouve que $Q'(\lambda,\zeta)=(\zeta^2\deleted{\unsure{+}}-\lambda^2\delta)/4$ est prop. à $Q$, et pour voir que le facteur de prop. est $1$, il suffit de faire \deleted{$y=(\lambda,\zeta)=T$}
% manuscript line(s) 8
\[
y=T,\qquad\text{i.e.}\quad y=(2,0)
\]
% manuscript line(s) 9
on a alors
% manuscript line(s) 9
\[
\deleted{Q(T)=T^2-2}
\]
% manuscript line(s) 10–11
\[
\begin{aligned}Q(y)&=y^2-Ty\psi(y)+N\psi(y)^2=T^2-2T^2+4N={}\\ &4N-T^2=\deleted{\unsure{+}}-\delta\end{aligned}
\]
% manuscript line(s) 12
et $Q'(y)=(0+4\delta)/\deleted{\unsure{8}}4=\unsure{-|+}\delta$ \qquad OK.
% manuscript line(s) 13
\par\centerline{\rule{1.1cm}{.3pt}}

% manuscript line(s) 14
Comparaison des torseurs sous $\mu_{2}$ et sous $(\mathbb Z/2\mathbb Z)_X$.
% manuscript line(s) 15
L'hom can.
% manuscript line(s) 16
\[
(\mathbb Z/2\mathbb Z)_X\longrightarrow\mu_{2\,X}
\]
% manuscript line(s) 17
définit, par torseur associé, un $\otimes$-foncteur
% manuscript line(s) 18–21
\diagram{\begin{tikzcd}[ampersand replacement=\&,column sep=small,row sep=large] \underline{\mathrm{Tors}}\ (\mathbb Z/2\mathbb Z)\arrow[r,"\mathrm{Ass}","\deleted{\unsure{\varphi}}"']\arrow[d,"\unsure{\wr}","\deleted{\unsure{\mathcal A}}"']\&\mathrm{Tors}(\mu_2\deleted{\unsure{,X}})\arrow[r,"\mathrm{Oub}"]\arrow[d,"\unsure{\alpha|\sim|\wr}"]\&{\big[\mathrm{Rev\,quad}(X)\big]}\arrow[d]\\ {\substack{\underline{\text{triples}}\ (L,\underset{\substack{\in\\\Gamma L^{\otimes2}}}{\delta},\underset{\substack{\in\\\Gamma L_0}}{T_0})\\(\underline{\delta\ \mathrm{inv.}})}}\arrow[r,"{\substack{\mathrm{Ass}'\\?}}"]\&{\substack{\underline{\text{couples}}\ (L,\underset{\substack{\in\\\Gamma L^{\otimes2}}}{\Theta})\\(\underline{\Theta\ \mathrm{inv}})}}\arrow[r,"{\substack{\mathrm{Oub}'\\?}}"]\&{\Big[\,\substack{\underline{\mathrm{Triples}}\ (L,\underset{\substack{\in\\\Gamma L^{\otimes2}}}{\delta},\underset{\substack{\in\\\Gamma L_0}}{T_0})\\(\delta\text{ pas néc\textsuperscript{t} inv.})}}\end{tikzcd}}
% manuscript line(s) 22–24
Les flèches verticales sont des équiv. si $2$ rég. dans $X$ (et celles \inserted{du milieu sans condition sur $X$}) Il faut déterminer les deux foncteurs horizontaux, notés ``?''. Le premier foncteur est sûrement
% manuscript line(s) 25
\[
\mathrm{Ass}'(L,\delta,T_0)\unsure{\longrightarrow|=}(L,\delta)
\]
% archive sheet number: 53
\scanpage{55}

% manuscript line(s) 1
Mais le deuxième serait plutôt
% manuscript line(s) 2
\[
\mathrm{Oub}'(L,\Theta)=(L,4\Theta,0)
\]
% manuscript line(s) 2–4
Ce foncteur ne commute \uline{pas} \inserted{néc\textsuperscript{t}} aux \deleted{\unsure{$\otimes$}}\unsure{$\wedge$|$\otimes$} à is près \inserted{(\deleted{\unsure{pas pour|N.B. pas}} $2$ inv. sur $X$,)}), et on en déduit que le foncteur
% manuscript line(s) 2–6
% margin note written sideways
% manuscript line(s) 2–6
\margin{car $(L\otimes L',16\Theta\Theta')$ n'est pas néc. isom. à $(L\otimes L',4\Theta\Theta')$, \deleted{\unsure{puisque pour avoir|p. ex.} $16\Theta\Theta'=0,\ 4\Theta\Theta'=0$} (p.ex. sur \deleted{\unsure{$Z$|$X$}} $\mathrm{Spec}\ \mathbb Z$, $L=L'=\underline O_X$, $\Theta\unsure{=}\Theta'=1$ $\ldots$}
% manuscript line(s) 5
\[
\mathrm{Tors}(\mu_{2/X})\unsure{\mathbin{-}|\to}\underline{\mathrm{Rev.quad}}(X)
\]
% manuscript line(s) 6
ne commute pas aux compositions, contrairement à
% manuscript line(s) 7–8
\diagram{\begin{tikzcd}[ampersand replacement=\&] \underline{\mathrm{Tors}}((\mathbb Z/2\mathbb Z)_X)\arrow[r]\arrow[d]\&\underline{\mathrm{Tors}}(\mu_{2X})\\ \underline{\mathrm{Rev.quad}}(X)\&\end{tikzcd}}

% manuscript line(s) 9
On aura en effet
% manuscript line(s) 10
\[
(L,\delta,T_0)\wedge(L',\delta',T'_0)=(L\otimes L',\delta\cdot\delta',T_0\cdot T'_0)
\]
% manuscript line(s) 11–12
\deleted{\unsure{tant|les}} des \uline{Triples} \uline{étales} $(L,\delta,T_0)$ que des \uline{Triples} (\unsure{quasi|quelconques} étales) $(L,\delta,T_0)$, et
% manuscript line(s) 13
\[
(L,\Theta)\wedge(L',\Theta')=(L\otimes L',\Theta\Theta')
\]
% manuscript line(s) 14
\par\centerline{\rule{1.2cm}{.3pt}}
% manuscript line(s) 14–19
% slanted
% manuscript line(s) 14–19
\margin{Ce qui suit (jusqu'à la page suivante) doit être repris de façon plus fine, en remplaçant \uline{Rev.quad!}$(X)$ par la cat. des $\mu_2$-rev. quadratiques (ou en \unsure{seulement|renonçant|se convaincant} \unsure{algèbres} $\ldots$)}

% manuscript line(s) 15–17
$\underline{\mathrm{Revquad}}(X)$ est une catégorie monoïdale (comm.) par un produit contracté qui prolonge celui sur les $(\mathbb Z/2\mathbb Z)_X$-torseurs

% manuscript line(s) 18–22
$\underline{\mathrm{Revquad}}^{!}(X)$, sous-catégorie pleine formée des rév. quadratiques \inserted{$X'/X$} tels que $T_0=0$ i.e. la trace sur $X'_0/X_0$ $(X_0=V(2.1_X))$ est id. nulle, i.e. $N_{X'_0/X_0}$ est additive i.e. $\varphi_{N_{X'_0/X_0}}$ (forme bil. associée) est id. nulle.

% manuscript line(s) 23
On définit une catégorie monoïdale mixte
% archive sheet number: 54
\scanpage{56}
% manuscript line(s) 1–4
des couples $(n,\unsure{\mathscr X|\mathcal X|\mathfrak X})$ où $n\in\mathbb Z/2\mathbb Z$ est une ``parité'', $\mathscr X/X$ rev. quadratique \uline{avec} \quad $X'\in\underline{\mathrm{Rev.quad}}^{!}$ \uline{si $n=1$}

% manuscript line(s) 5
On pose
% manuscript line(s) 6–7
\[
(n,\deleted{\unsure{X'}}\mathscr X)\wedge(n',\mathscr X')=\begin{cases}(n+n',\mathscr X\wedge\mathscr X')&\text{si $n$ ou $n'$ pair}\\ (n+n',\mathscr X\overset{!}{\wedge}\mathscr X')&\text{si $n$ \uline{et} $n'$ impairs}\end{cases}
\]
% manuscript line(s) 8–10
où $\mathscr X\overset{!}{\wedge}\mathscr X'$ est un produit contracté modifié, défini si $\mathscr X,\mathscr X'\in\underline{\mathrm{Rev.quad}}^!(X)$, et défini (si $2$ rég. sur $X$) par la condition
% manuscript line(s) 11
\[
(L,\delta,0)\overset{!}{\wedge}(L',\delta',0)=(L\otimes L',\tfrac14\delta\delta',0)
\]
% manuscript line(s) 12–13
où $\tfrac14\delta\delta'$ a un sens intrinsèque, car loc.\textsuperscript{t} on a
% manuscript line(s) 13–14
\[
\begin{aligned}\delta&=4\Delta &&\Delta\ \text{définie mod }h\ \text{avec }4h=0\\ \delta'&=4\Delta' &&\Delta'\ \text{\rule{1.2cm}{.3pt} }h'\ \text{\rule{1.2cm}{.3pt} }4h'=0\end{aligned}
\]
% manuscript line(s) 15
\[
\delta\delta'=16\Delta\Delta'\qquad\tfrac14\delta\delta'\overset{\mathrm{def}}=4\Delta\Delta'
\]
% manuscript line(s) 15–16
qui ne change pas \unsure{--} si $\Delta\longmapsto\Delta+h\quad(4h=0$ \quad $\Delta'\longmapsto\Delta'+h'\quad(4h'=0)$

% manuscript line(s) 17–19
\leavevmode\llap{\rule[-2\baselineskip]{.3pt}{3\baselineskip}\hspace{2pt}\rule[-2\baselineskip]{.3pt}{3\baselineskip}\quad}Il faut lui donner un sens indépendamment de la condition $2$ régulier sur $X$, exercice catégorique universel $\ldots$

% manuscript line(s) 20–24
Cette opération est \uline{commutative} \inserted{(à isom. can. près)} et de plus associative, les isom de comm. et d'assoc. étant les habituels sur $L\otimes L',\deleted{\unsure{\simeq}}L'\otimes L$, et $(L\otimes L')\otimes L''\simeq L\otimes(L'\otimes{\unsure{\underline L|L}}'')$, et respectant automatiquement les pondérations $\delta,T_0$ etc.
% archive sheet number: 55
\scanpage{57}

% manuscript line(s) 1
\uline{Rev quadratiques en car. $2$}
% manuscript line(s) 2
On a
% manuscript line(s) 2
\[
L\rightrightarrows L^{\otimes2}\qquad\begin{cases}u\longrightarrow u^2\\u\longrightarrow Tu\end{cases}\qquad T\in\Gamma L\ \text{la section trace}
\]
% manuscript line(s) 3
(\uline{NB} ce sont des hom. additifs, mais $u\mapsto u^2$ n'est pas linéaire)
% manuscript line(s) 4
d'où
% manuscript line(s) 4
\[
L\xrightarrow{f}L^{\otimes2}:\quad f(u)=u^2+Tu
\]
% manuscript line(s) 5
et $\check V(L)\underset f\longrightarrow\check V(L^{\otimes2})$ morphisme qui est étale ssi $T$ est inv.
% manuscript line(s) 6
\deleted{On peut considérer} et alors $T:\underline O_X\xrightarrow{\sim}L$, et $f$ peut s'interpréter comme $\underline O_X\longrightarrow\underline O_X:u\longmapsto u^2+u$

% manuscript line(s) 7–8
Considérons \deleted{\unsure{$x'$ et $x''$|$x'$, $x''$}} les deux sections éventuelles $x',x''$ de $X'\subset P_1\subset E$, on a
% manuscript line(s) 9
\[
x''\deleted{\unsure{+}}\unsure{-|+}x'=x''+x'=T\qquad\text{i.e}\qquad x''=x'+T
\]
% manuscript line(s) 10
d'où
% manuscript line(s) 10
\[
{x''}^{(2)}={x'}^{(2)}+T^2\qquad\text{sections de }P_1\mathbin{*}_L L^2\simeq P_1^{(2)}
\]
% manuscript line(s) 11–12
\diagram{\begin{tikzcd}[ampersand replacement=\&,column sep=small]0\arrow[r]\&L\arrow[r]\arrow[d,"{\xi\mapsto\xi^2}"]\&E\arrow[r,"\psi"]\arrow[d,"{x\mapsto x^{(2)}}"]\&\underline O_X\arrow[r]\arrow[d,"{\lambda\mapsto\lambda^2}"]\&0\\0\arrow[r]\&L^{\otimes2}\arrow[r]\&E^{(2)}\arrow[r,"\psi^{(2)}"]\&\underline O_X\arrow[r]\&0\end{tikzcd}}
% manuscript line(s) 12
\margin{$P_1^{(2)}=\psi^{(2)-1}(\{1\})$}
% manuscript line(s) 13
\[
\underset{\mathrm{Sym}^2(E)}{\deleted{\unsure{\mathrm{Sym}^2(E)}}}\big/\left\{(\,x\deleted{\unsure{-}}+y)^{(2)}-\underbrace{x^{(2)}-y^{(2)}}_{xy}\right\}\deleted{\unsure{x^{(2)}=y^{(2)}}}
\]
% manuscript line(s) 14–16
\[
\begin{aligned}\unsure{w|u}\in\Gamma(P_1^{(2)}\wedge T\cdot P_1)&\Longleftrightarrow {\unsure{w|u}}': P_1^{(2)}\simeq TP_1\quad\text{iso de }L^{\otimes2}\text{-torseurs}\\ &\Longleftrightarrow P_1\xrightarrow{u''}P_1^{(2)}\quad\text{compatible avec}\\ &\qquad L\longrightarrow L^{\otimes2}\\ &\qquad u\longmapsto T\cdot u\end{aligned}
\]

% manuscript line(s) 17–20
\uline{NB} $f:\check V(L)\mathbin{-}\check V(L^{\otimes2})$ est fid plat (donc épi fppf) et son noyau est un sous-groupe \inserted{fini} quadratique de $\check V(L)$, soit, \deleted{\unsure{$K$}} $G_T$, étale ssi $T$ est inv. (son ens de ramification est donc $V(T)$, schéma des zéros de $T$)

% manuscript line(s) 21–22
Soit $x\in\deleted{\unsure{\Gamma(P_1)}}\inserted{\Gamma X'}$, i.e. $x\in\Gamma(P_1)$, avec
% manuscript line(s) 22
\[
Q(x)=0\qquad x^2\underset{\underbrace{\text{ou }+}}{-}Tx+N=0
\]
% manuscript line(s) 23
Écrivons $Q(x+u)=0$ \inserted{(où $u\in\Gamma\underline L$)}, \quad $\underset{/}{x^2}+u^2+\underset{/}{Tx}+Tu+\underset{/}{N}=0$, i.e. $u^2+Tu=0$
% archive sheet number: 56
\scanpage{58}

% manuscript line(s) 1
Donc $X'/X$ est un \uline{torseur} sous $G_T$ ! Et \unsure{ce|se|si|2}

% manuscript line(s) 2–6
\uline{Th} On a une équivalence de catégories \inserted{(groupoïdes)} fibrées sur $\mathrm{Sch}/\mathbb F_2$ (plus gén. sur les topos loc\textsuperscript{t} annelés en car. $2$), \unsure{entre} la catégorie des rev. quadratiques relatifs (sur $X\in\mathrm{Ob}\ \mathrm{Sch}/\mathbb F_2$), et les données suivantes sur un tel $X$
% manuscript line(s) 7–10
\[
(L,\inserted{P_1,}T,\alpha),\quad\text{où}\quad\left\{\begin{aligned}&L\text{ module inv. sur }X,\ P_1\ L\text{-tors.}\\ &T\text{ section de }L\ \text{(module des translations de }P_1\text{)}\\ &\alpha\text{ section de}\\ &\quad P_1\wedge_{\underline L}(L^{\otimes2},f_T:L\unsure{\mathbin{-}|\to}L^{\otimes2})\end{aligned}\right.
\]
% manuscript line(s) 7
\margin{(i.e. $P_1$ un fibré en droites \uline{affines} sur $X$)}
% manuscript line(s) 11
où $f_T:L\mathbin{-}L^{\otimes2}$ donné par $f_T(u)=u^2+T\cdot u$
% manuscript line(s) 12–13
(c'est un \unsure{homo|hom.} additif $W(L)\mathbin{-}W(L^{\otimes2})$, dont le noyau est un \inserted{schéma en} \unsure{groupe|groupes} \inserted{fini} quadratique $G_T$ sur $X$).
% manuscript line(s) 14
Au système $\Sigma=(L,P_1,T,\alpha)$ est associé le $X$-schéma
% manuscript line(s) 15
\[
X'_\Sigma=f_{T,P_1}\unsure{\zeta^{-1}|{}^{-1}}(V(\alpha))
\]
% manuscript line(s) 16
(torseur sous $G_T$), qui est \inserted{fini} quadratique sur $X$.

% manuscript line(s) 17
Plus de doute, c'est lui !
% manuscript line(s) 17–22
\margin{\uline{NB} $u:P_1\longrightarrow P_1^{(2)}$ comp. avec $u\mapsto Tu$ est donné par $u:x\longmapsto Tx+N$, ou $\underline{u(x)=x^2+Q(x)}$ en effet $(Q(x)=x^2+Tx+N)$, donc \deleted{\unsure{on aura}} \deleted{\unsure{$X'$}} l'équation $Q(x)=0$ qui car. $X'/X$ équivaut à $\boxed{u(x)=x^2}$ \quad $P_1\mathrel{\substack{\xrightarrow{u}\\[-3pt]\xrightarrow[\ x\mapsto x^2\ ]{\phantom{u}}}}P_1^{(2)}$ \deleted{\unsure{$f$}}}

% manuscript line(s) 18–19
Comment s'explicite le produit $\wedge$ des rev. quadratiques ?
% manuscript line(s) 20–21
\[
\begin{aligned}(L,T,P_1,\alpha)\wedge(L',T',P'_1,\alpha')&\overset?={}\\ &(L\otimes L',TT',?,?)\end{aligned}
\]
% manuscript line(s) 22–23
\unsure{Ainsi|Alors} :
% manuscript line(s) 22–23
\[
\begin{aligned}G_T\xrightarrow{f_{T'}}G_{TT'}&\quad\text{induit par } f_{T'}:L\mathbin{-}L\otimes L',\quad u\longmapsto uT'\\G_{T'}\mathbin{-}G_{TT'}&\quad\rule{1cm}{.3pt}\ f_T:L'\mathbin{-}L\otimes L',\quad u'\mathbin{-}Tu'\end{aligned}
\]
% manuscript line(s) 24–26
d'où $G_T\times G_{T'}\longrightarrow G_{TT'}$ hom de groupes d'où, pour des torseurs $\unsure{\mathscr X|\mathcal X|\mathfrak X}$ sous $G_T$, $\mathscr X'$ sous $G_{T'}$, un torseur $\mathscr X\wedge\mathscr X'$ sous $G_{TT'}$ !
% archive sheet number: 57
\scanpage{59}

% manuscript line(s) 1–3
\deleted{Trois} \inserted{Quatre} catégories fibrées sur $(\mathrm{Sch}_{\unsure{?}})$ (plus gén. sur la catégorie des topos loc. annelés $\ldots$), canoniquement équivalentes
% manuscript line(s) 4–11
\diagram{\begin{tikzcd}[ampersand replacement=\&,column sep=small,row sep=small] \text{\textcircled{1}}\arrow[d,leftrightarrow]\&\parbox{.78\linewidth}{Catégorie des rev. quadratiques $X'/X$ de $X\in\mathrm{Ob}\,\mathrm{Sch}$,}\\ \text{\textcircled{$\deleted{\unsure{2}}1'$}}\arrow[d]\&\parbox{.78\linewidth}{Catégorie des alg \inserted{$\mathcal A$} commutatives loc libres de rang $2$ sur $\underline O_X$}\\ \text{\textcircled{2}}\&\parbox{.78\linewidth}{Catégorie des triples $(\mathcal A,N,e)$, où $(\mathcal A,N)$ est un module quadratique loc. libre de rang $2$ sur $\underline O_X$, et où $e$ est une section de $\mathcal A$ telle que \(\displaystyle N(e)=1\)}\\ \text{\textcircled{3}}\arrow[uuu,bend left=35]\&\parbox{.78\linewidth}{Catégorie des triples $(L,P,Q)$, où $L$ est un $\underline O_X$-Module inv. \inserted{(identifié au schéma en droites homogène corresp)}, $P$ un $L$-torseur, et $Q$ une application polynomiale}\end{tikzcd}}
% manuscript line(s) 12
\[
P\longrightarrow L^{\otimes2}
\]
% manuscript line(s) 13–14
dont la partie homogène de degré $2$ soit $u\unsure{\mathbin{-}|\mapsto}u^2$, i.e. on suppose que
% manuscript line(s) 15
\[
Q(x+u)=Q(x)+u^2+\mathrm{Lin}_u(x)
\]
% manuscript line(s) 16–18
où $\mathrm{Lin}_u:P_{S}\mathbin{-}L_S^{\otimes2}$ est une \unsure{application|appl.} \deleted{\unsure{linéaire|lin.}} \uline{affine} (quand $u$ est une section loc.\textsuperscript{le} \inserted{fppf, disons} sur $S/X$)

% manuscript line(s) 19–22
$\updownarrow$\quad\textcircled{$3'$} Catégorie des triples $(0\to L\overset{i}\to E\overset{\psi}\to\underline O_X\to0,Q)$ d'une suite exacte $0\mathbin{-}L\mathbin{-}E\mathbin{-}\underline O_X\mathbin{-}0$, avec $L$ inv., et d'une forme quadratique $Q$ sur $E$ à val. dans $L^{\otimes2}$, dont la restriction à $L$ soit $u\mathbin{-}u^{\otimes2}$.

% manuscript line(s) 23–25
On va transformer \textcircled{$3'$}, en \deleted{\unsure{pl}} considérant $L^{\otimes2}$ comme plongé dans $\mathrm{Sym}^2(E)$ ($\simeq\mathrm{Quad}(\mathcal A)$, où $\mathcal A\overset{\mathrm{def}}=\check E$), lequel a une suite de composition
% archive sheet number: 58
\scanpage{60}
% manuscript line(s) 1
\[
\underbrace{0\subset L^{\otimes2}}_{L^{\otimes2}}\ \underbrace{\subset L\otimes E}_{L^{\deleted{\unsure{\otimes2}}}}\ \underbrace{\subset\mathrm{Sym}^2(E)}_{\underline O_X}
\]

% manuscript line(s) 2
Considérons \deleted{\unsure{écrivons}}
% manuscript line(s) 3
\[
Q(x)-x^2:\quad\unsure{P|E}\mathbin{-}\mathrm{Sym}^2(E)
\]
% manuscript line(s) 4–6
fonction \deleted{quadratique, sa partie homogène est} \inserted{polynomiale de degré \deleted{\unsure{possiblement}} $\leq2$, dont la} partie homogène \deleted{\unsure{est sa}} quadratique
% manuscript line(s) 6
\[
\deleted{Q(x+u)=}\qquad L\mathbin{-}L^{\otimes2}
\]
% manuscript line(s) 7
est nulle. Donc c'est une application \uline{affine}

% manuscript line(s) 8–12
\deleted{soit $Q(x)=x^2+\mathrm{Af}(x)$ \quad $\mathrm{Af}:P\mathbin{-}\deleted{\unsure{E}}\mathrm{Sym}^2 E$ application \uline{affine}. Son terme linéaire est un hom $L\longrightarrow\deleted{\unsure{E}}\mathrm{Sym}^2(E)$, \inserted{donc} défini par un \unsure{tel} élément $T$ de $L^\vee\otimes\mathrm{Sym}^2 E\supset E$}

% manuscript line(s) 13–15
Utilisant un scindage de \unsure{$N$ cer. $E$|l'ext. $E$|$N$ et $E$} par une section $e$ de $P_1$, i.e. $e\in\Gamma E$, $\psi(e)=1$, on aura, posant
% manuscript line(s) 15
\[
\begin{cases}x=e+u\\ u=x-e\end{cases},\quad\inserted{u\in\Gamma\underline L,}\quad\text{et écrivant}
\]
% manuscript line(s) 16
\[
Q(e+u)=u^2+bu+c\qquad\begin{aligned}u&\in\Gamma L\\ b&\in\Gamma\unsure{L'|L}\\ c&\in\Gamma L^{\otimes2}\end{aligned}
\]
% manuscript line(s) 17–18
\[
\begin{aligned}Q(\underbrace{e+u}_{x})-\deleted{\unsure{Q}}(\underbrace{e+u}_{x})^2&=bu+c-2eu-e^2=\underbrace{u}_{x-e}(b-2e)+c-e^2\\ &=\underline x(\unsure{+}b-2e)\deleted{\unsure{-e(b-e)+c-e^2}}+(e^2+c-be)\end{aligned}
\]
% manuscript line(s) 19–21
\[
\text{i.e.}\quad\left\{\begin{aligned}Q(x)&=x^2\deleted{\unsure{+}}\unsure{-|+}xT+N\\[-2pt] &\text{où}\quad\boxed{T=2e-b\in E,\quad N=e^2+\underset{\substack{\in\\\Gamma L^2}}{c}-\underset{\substack{\in\\\Gamma(L\otimes E)\subset\Gamma(\mathrm{Sym}^2E)}}{be}\in\mathrm{Sym}^2(E)}\\ &\text{avec}\quad\boxed{\psi(T)=2,\varepsilon(N)=1}\end{aligned}\right.
\]
% manuscript line(s) 21
\diagram{\begin{tikzcd}[ampersand replacement=\&,column sep=small]0\arrow[r,dash]\&L\otimes E\arrow[r]\&\mathrm{Sym}^2(E)\arrow[r,dash,"\varepsilon"]\&\underline O_X\arrow[r,dash]\&0\end{tikzcd}}

% manuscript line(s) 22–26
\deleted{\illegible{2}} On vérifie que $T$, donc aussi $N$, ne dépend pas du scindage choisi. [$T$ est élém\textsuperscript{t} d'un torseur \inserted{$P_2=\psi^{-1}(2)$} sous $L^{\deleted{\unsure{\otimes}}}$, \deleted{\unsure{donc|et}} et $N$ (pour $T$ choisi), \deleted{que} la condition que $N$ soit à valeurs dans $L^{\otimes2}$, est élém\textsuperscript{t} d'un torseur sous $L^{\otimes2}\subset\mathrm{Sym}^2(E)$, donc \inserted{le choix de} $Q\deleted{\unsure{+2}}\longleftrightarrow(T,N)$ est \unsure{celui}
% archive sheet number: 59
\scanpage{61}
% manuscript line(s) 1–4
d'un élément d'un torseur sous $L^{\otimes2}\otimes\check E$ \inserted{$=L^{\otimes2}\otimes\mathcal A$}, qui \deleted{\unsure{est}} peut être considéré comme l'espace des applications \uline{affines} de $P_1$ dans $L^{\otimes2}$, \inserted{i.e. des applications linéaires de $E$ dans $L^{\otimes2}$, soit $\check E\otimes L^{\otimes2}$} \deleted{\unsure{donc c'est}} et on a
% manuscript line(s) 5
\[
0\unsure{\mathbin{-}|\to}L^{\otimes2}\longrightarrow L^{\otimes2}\otimes\check E\mathbin{-}L\longrightarrow0
\]
% manuscript line(s) 5–7
\margin{(car $0\mathbin{-}\underline O_X\mathbin{-}\check E\mathbin{-}\check L\longrightarrow0$) et on multiplie par $L^{\otimes2}$}
% manuscript line(s) 6–8
$\underset{L\otimes E}{\unsure{\shortparallel|?|\simeq}}$ \uline{can} (pour tt module loc. lib. de rg $2$, \unsure{$F$|$E$} avec $L\simeq\det E$) i.e.
% manuscript line(s) 8
\[
E\simeq\check E\otimes L\ \big(\simeq\underline{\mathrm{Hom}}(E,\underset{\substack{\shortparallel\\\det E}}L)\big)
\]

% manuscript line(s) 9
L'application affine
% manuscript line(s) 10
\[
P_1\longrightarrow L^{\otimes2}
\]
% manuscript line(s) 11
définie par un $\unsure{\zeta|\xi}\in\underset{\substack{\cap\\E}}L\cdot E\simeq L\otimes E\subset\mathrm{Sym}^2E$
% manuscript line(s) 12
est définie par
% manuscript line(s) 13
\[
x\longmapsto x\deleted{\unsure{\otimes}}\wedge\zeta
\]
% manuscript line(s) 12–15
\margin{$E\otimes E\to L$ \quad $(x\otimes y)\mapsto x\wedge y$ \quad $\underbrace{L\otimes E}_{\ni\,\zeta}\otimes\underset{\substack{\ni\\x}}E\mathbin{-}L\otimes L=L^{\otimes2}$ \quad $(\zeta,x)\mathbin{-}\zeta\wedge x=-x\wedge\zeta$}
% manuscript line(s) 16
$\ldots\,.$]

% manuscript line(s) 17–18
Inversement, si (l'ext. $E$ de $\underline O_X$ par $L$ étant donnée, sans $Q$ encore) si on se donne
% manuscript line(s) 19
\[
T\in E,\ N\in\mathrm{Sym}^2(E),
\]
% manuscript line(s) 20–21
quand correspondent-ils à une $Q$ comme dans 3) ou 3$'$) ? Il faut, rappelons le
% manuscript line(s) 22
\[
\psi_1(T)=2,\qquad\unsure{\psi_2|\varphi_2}(N)=1
\]
% manuscript line(s) 22–23
\margin{i.e. (en termes de la ``section unité'' $e$ de $\mathcal A=\check E$ : $T(e)=2,N(e)=1$)}
% manuscript line(s) 24
plus une condition qui équivaudrait à
% manuscript line(s) 25
\[
N(x+\deleted{\unsure{\lambda}}e)=N(x)+\deleted{\unsure{\lambda}}T(x)+1^{\deleted{\unsure{2}}}
\]
% manuscript line(s) 25
\margin{(\inserted{on aura $N(x+\lambda e)=N(x)+\lambda T(x)+\lambda^2$})}
% manuscript line(s) 26
i.e. $T(x)=\varphi_N(x,e)$

% manuscript line(s) 27
En tous cas on doit avoir : \uline{$N$ détermine $T$}
% archive sheet number: 60
\scanpage{62}

% manuscript line(s) 1
Je vois, du côté $E$, la condition que
% manuscript line(s) 2
\[
Q_{\deleted{\unsure{x}}T,N}(x)\overset{\mathrm{def}}=x^2\deleted{\unsure{+}}\unsure{-|+}xT+N\qquad(x\in\Gamma P_1)
\]
% manuscript line(s) 3
à valeurs dans $L^{\otimes2}$, ou \unsure{mieux|encore}
% manuscript line(s) 4
\[
Q_{T,N}(x)=x^2-\psi(x)xT+\psi(x)^2N\qquad(x\in\Gamma\underline E)
\]
% manuscript line(s) 5–6
à valeurs dans $L^{\otimes2}$. Faisant $x=T$, on trouve la condition familière (compte tenu de $\psi(T)=2$)
% manuscript line(s) 7
\[
4N-T^2\in\Gamma L^{\otimes2}\quad\subset\Gamma\mathrm{Sym}^2 E
\]
% manuscript line(s) 8–10
mais elle ne suffit pas. En termes d'un scindage $E\simeq L\oplus\underline O$ par choix d'un $x_0\in\Gamma P_1$, on aura
% manuscript line(s) 10
\[
\mathrm{Sym}^2E\simeq\underline L^2\oplus\underset{\substack{\shortparallel\\x_0\cdot L}}{\underline L}\oplus\underset{\substack{\phantom{\shortparallel}\\=x_0^2\underline O_X}}{\underline O}
\]
% manuscript line(s) 11–13
\[
\begin{aligned}&\deleted{\unsure{y}}\quad x_0=(0,1)\deleted{\unsure{x}}\\ &T=(t,2)=2x_0+\underset{(t,0)}t,\quad t\in\Gamma L,\\ &N=\underset{\substack{\in\\\Gamma L^{\otimes2}}}c-x_0\underset{\substack{\in\\\Gamma L}}b+x_0^2\quad(\text{car }\psi_2(N)=1)\end{aligned}\qquad\begin{cases}t\in\Gamma L\\ b\in\Gamma L\\c\in\Gamma L^{\otimes2}\end{cases}
\]
% manuscript line(s) 14
Soit alors $x=(u,1)\inserted{=u+x_0}$, $u\in L$, on aura
% manuscript line(s) 15–17
\[
\begin{aligned}Q_{T,N}(x_0+u)=Q_{T,N}(x)&=(x_0+u)^2-(x_0+\underset{\substack{\in\\\Gamma L}}u)(2x_0+\underset{\substack{\in\\\Gamma L}}t)+\underset{\substack{\in\\\Gamma L^{\unsure{2}}}}c-\underset{\substack{\in\\x_0\otimes L}}{bx_0}+x_0^2\\ &\deleted{\unsure{=2x_0u+x_0^2+u^2-(x_0+u)t}}\\ &=\underset{\substack{\in\\L^{\otimes2}}}{u^2}-\underset{\substack{\in\\L^{\otimes2}}}{ut}+\deleted{\unsure{x_0^2+}}(-bx_0+c-x_0t)\end{aligned}
\]
% manuscript line(s) 18
et la condition est que
% manuscript line(s) 19–20
\[
\begin{aligned}&\deleted{x_0^2-x_0(x_0}\\&\underbrace{x_0(\deleted{\unsure{x_0+}}-b-t)}_{\in\Gamma(E\otimes L)}+\underset{\substack{\in\\\Gamma(L^{\otimes2})\subset\Gamma(E\otimes L)}}c\in\Gamma L^{\otimes2}\end{aligned}
\]
% manuscript line(s) 21
Cela signifie $\boxed{t=-b}$ i.e. la condition est \deleted{que}
% manuscript line(s) 21–22
% slanted
% manuscript line(s) 21–22
\margin{Ceci \deleted{signifie} \inserted{implique} que sur le \unsure{composant} du scindage de l'ext. $\mathcal A$ de $\underline O$ par $L$, $b$ est le \unsure{même} (le \unsure{terme en} \unsure{bon} au signe}
% manuscript line(s) 22
\[
\boxed{N=x_0^{\unsure{2}}-\underset{\substack{\in\\\Gamma L}}b x_0+\underset{\substack{\in\\\Gamma L^{\otimes2}}}c\quad\Longrightarrow\quad T=2x_0\deleted{\unsure{+}}-b}
\]
% manuscript line(s) 22–23
% slanted
% manuscript line(s) 22–23
\margin{d'où $\boxed{\unsure{\Theta|D|\delta}=T^2-4N=b^2-4c}$}
% manuscript line(s) 23
d'où $\boxed{Q(x_0+u)=u^2+\deleted{\unsure{t}}bu+c}$
% archive sheet number: 61
\scanpage{63}

% manuscript line(s) 1–2
Ainsi on trouve équivalence $\underset{\substack{\wr\\3)}}{\unsure{3')|3)|1)}}\simeq2)$ d'où \deleted{\unsure{la}} (via $3)\longrightarrow1)$) la flèche $2)\longrightarrow1)$.
% manuscript line(s) 3
\figure{\begin{tikzpicture}[x=1cm,y=1cm]\draw (1,.14) .. controls (1.3,-.03) and (2.1,.24) .. (2.8,.23);\draw (0,0).. controls (1.1,-.14) and (2.6,.1)..(3.1,.07);\draw (.4,-.29).. controls (.55,-.12) and (2.42,-.11)..(2.92,-.12);\draw (1.9,.07)--(2.12,-.12);\end{tikzpicture}}
% manuscript line(s) 4–6
% slanted
% manuscript line(s) 4–6
\margin{Produit contracté de rev. quadratiques}

% manuscript line(s) 4–6
On veut maintenant définir le foncteur ``\uline{produit contracté}'' sur la catégorie des rev. quadratiques. On veut
% manuscript line(s) 7–11
\[
\left\{\begin{aligned}\mathbb L(X_1\wedge X_2)&\simeq\mathbb L(X_1)\otimes\mathbb L(X_2)\\ &\hspace{-5em}\substack{\text{(module déterminantiel}\\\text{du rev. quadratique,}\\\text{noté }\underline L\text{ ci-dessus)}}\\ \underset{\substack{\in\\\Gamma\mathbb L(X_1\wedge X_2)}}{\delta_{X_1\wedge X_2}}&=\delta_{X_1}\cdot\delta_{X_2}\end{aligned}\right.
\]

% manuscript line(s) 12–16
De plus, on veut que tout scindage de $X_1$, \deleted{\unsure{qui}} défini par $x_1\in E_1$ avec $\psi_{X_1}(x_1)=1$, et tout scindage de $X_2$, défini par $x_2\in E_2$ avec $\psi_{X_2}(x_2)=1$, en définisse un de $X_1\wedge X_2$, i.e. on veut une application
% manuscript line(s) 17
$(x_1,x_2)\longrightarrow x_1\wedge x_2$ :
% manuscript line(s) 17–18
\diagram{\begin{tikzcd}[ampersand replacement=\&] P_{X_1}\&{}\times{}\&P_{X_2}\arrow[r]\&P_{X_1\wedge X_2}\\ X_1\arrow[u,hook]\&{}\times{}\&X_2\arrow[u,hook]\arrow[r,dash,"\text{qui induise}"']\&X_1\wedge X_2\arrow[u,hook]\end{tikzcd}}

% manuscript line(s) 19–21
Si $X_1,X_2$ munis des scindages $x_1,x_2$, et si les modules \inserted{inv.} \unsure{discriminants|discriminant} sont $L_1,L_2$, alors les structures sont définies par
% manuscript line(s) 22–23
\[
\begin{aligned}b_1&\in\Gamma L_1&b_2&\in\Gamma L_1^{\otimes2}&b_1&=2x_1-T_1\\ \unsure{b_2|b_1|c_1}&\in\Gamma L_2&c_2&\in\Gamma L_2^{\otimes2}&b_2&=2x_2-T_2\end{aligned}
\]
% archive sheet number: 62
\scanpage{64}
% manuscript line(s) 1
et on doit avoir alors (sauf erreur)
% manuscript line(s) 2
\[
\boxed{b\overset{\mathrm{def}}=b_{X_1\wedge X_2}=b_1b_2}.
\]
% manuscript line(s) 3–5
(on se demande ce que ça serait d'autre, au signe près $\ldots$), d'où pratiquement aussi $c$, via
% manuscript line(s) 6
\[
4c=b^2-\delta=\underbrace{(b_1^{\deleted{\unsure{2}}}b_2)^2}_{\substack{\deleted{\unsure{=b_1}}b_1^2b_2^2\\\shortparallel\qquad\shortparallel\\(\delta_1+4c_1)(\delta_2+4c_2)}}-\delta_1\delta_2=4(c_1\delta_2+c_2\delta_1+4c_1c_2)
\]
% manuscript line(s) 7
et on va faire
% manuscript line(s) 8
\[
\boxed{c=c_1\underset{\substack{\shortparallel\\b_2^2-4c_2}}{\delta_2}+c_2\underset{\substack{\shortparallel\\b_1^2-4c_1}}{\delta_1}+4c_1c_1}\ =\underline{b_2^2c_1+b_1^2c_2-4c_1c_2}
\]

% manuscript line(s) 9
Mais quelle est la formule pour
% manuscript line(s) 10
\[
(x_1+\underset{\substack{\in\\\Gamma L_1}}{u_1})\wedge(x_2+\underset{\substack{\in\\\Gamma L_2}}{u_2})\overset{?}=x_1\wedge x_2+\underbrace{f_{x_1,x_2}(u_1,u_2)}_{?}
\]
% manuscript line(s) 11
où $f_{x_1,x_2}:\deleted{\unsure{\Gamma}}L_1\times L_2\longrightarrow L_1\otimes L_2$.
% manuscript line(s) 12
On a trois applications \deleted{\unsure{de} $L_1\times L_2$ dans $L_1\otimes L_2$} \inserted{\unsure{can.} à val. \unsure{ds|dans} $L_1\otimes L_2$}
% manuscript line(s) 13
\[
\deleted{\unsure{f_1,f_2\text{ et }g:L_1\times L_2\to L_1\otimes L_2}}\qquad\deleted{\unsure{\xi_1\mapsto \xi_1\otimes}}
\]
% manuscript line(s) 14–16
\[
\begin{aligned}f_{x_2}:L_1&\unsure{\mathbin{-}|\to}L_1\otimes L_2&u_1&\mathbin{-}u_1b_2\\ f_{x_1}:L_2&\mathbin{-}L_1\otimes L_2&u_2&\mathbin{-}u_2b_1\\ \unsure{f\ell_{x_1,x_2}|f_{x_1,x_2}}:L_1\times L_2&\longrightarrow L_1\otimes L_2.&(\unsure{u_1|x_1},u_2)&\mathbin{-}u_1\otimes u_2\end{aligned}
\]
% manuscript line(s) 17
Je devine qu'on aura
% manuscript line(s) 18
\[
\boxed{(x_1+u_1)\wedge(x_2+u_2)=x_1\wedge x_2+\big(\underbrace{u_1b_2+u_2b_1+2u_1u_2}_{\in\Gamma(L_1\otimes L_2)}\big)}
\]
% archive sheet number: 63
\scanpage{65}

% manuscript line(s) 1–2
Vérifions que (pour $x_1,x_2$ fixés) \inserted{(i.e. pour $(L_1,b_1,c_1)$ et $(L_2,b_2,c_2)$ fixés} cette application
% manuscript line(s) 3
\[
\deleted{\unsure{L_1}}P_1\times\deleted{\unsure{L_2}}P_2\longrightarrow P\qquad\text{i.e.}\quad L_1\deleted{\unsure{\otimes}}\times L_2\longrightarrow L_1\otimes L_2
\]
% manuscript line(s) 4
\deleted{\unsure{identifie|applique}} applique bien $X_1\times X_2$ dans $X$, i.e.
% manuscript line(s) 5–8
\[
\begin{aligned}\bigg|\, &\overbrace{u_1^2+b_1u_1+c_1}^{Q_1(u_1)}=0\\ &\underbrace{u_2^2+b_2u_2+c_2}_{Q_2(u_2)}=0\\ \Longrightarrow\ &\underbrace{u^2+bu+c}_{Q(u)}=0\qquad\text{où}\quad\left\{\begin{aligned}u&=\overbrace{u_1b_2+u_2b_1}+\deleted{\unsure{+}}2u_1u_2\\ b&=b_1b_2\\ c&=b_2^2c_1+b_1^2c_2\deleted{\unsure{+}}-4c_1c_2\end{aligned}\right.\end{aligned}
\]
% manuscript line(s) 9
\[
\big(\overset{\frown}{\underset{\smile}{\overset{\frown\qquad\frown}{\underset{\smile\qquad\smile}{\underset{\cdot}{\overset{\cdot}{u_1b_2}}+\underset{\cdot}{\overset{\cdot}{u_2b_1}}+\underset{\cdot}{\overset{\cdot}{2u_1u_2}}}}}}\big)^2+\underbrace{b_1b_2\big(\overset{\cdot}{u_1b_2}+\overset{\cdot}{u_2b_1}+\deleted{\unsure{u_1u_2}}2\overset{\cdot}{u_1u_2}\big)}+\overset{\cdot}{b_2^2c_1}+\overset{\cdot}{b_1^2c_2}-\overset{\cdot}{4c_1c_2}=\unsure{0\,?}
\]
% manuscript line(s) 10
\deleted{\unsure{$b_2^2$}} Le premier mb s'écrit ($\forall u_1\in\Gamma L_1,\ u_2\in\Gamma L_2$
% manuscript line(s) 11–14
\[
\begin{aligned}&b_2^2(u_1^2+b_1u_1+c_1)+b_1^2(u_2^2+b_2u_2+c_2)+\underbrace{4u_1^2u_2^2}_{\unsure{4Q_{u_1}\cdot Q_{u_2}}}\\ &\qquad\deleted{\unsure{+4u_1u_2b_1b_2}}\quad\deleted{\unsure{+4b_1b_2u_1u_2}}\qquad\deleted{\unsure{Q(u_1)=b_1u_1+c_1}}\\ &+4u_1^2b_2u_2+4u_2^2b_1u_1-4c_1c_2\deleted{\unsure{+2}}\\ &\qquad\deleted{\unsure{+2u_1u_2b_1b_2|+4u_1u_2b_1b_2}}\end{aligned}
\]
% manuscript line(s) 15–18
\[
\begin{aligned}&=b_2^2Q_1(u_1)+b_1^2Q_2(u_2)+4\big(Q(u_1)-c_1\big)\big(Q(u_2)-c_2\big)\\ &\qquad -4b_1b_2u_1u_2\deleted{\unsure{-4c}}\quad\deleted{\unsure{+4b_1b_2u_1u_2|-4b_1b_2u_1u_2}}-4c_1c_2\\ &=b_2^2Q_1(u_1)+b_1^2Q_2(u_2)+4\big[Q(u_1)Q(u_2)-c_1Q(u_2)-c_2Q(u_1)\big]\\ &\phantom{{}={}}\deleted{\unsure{+4b_1b_2u_1u_2|-4b_1b_2u_1u_2}}\end{aligned}
\]
% manuscript line(s) 19–20
\[
\overset{u^2-bu+c=}{\boxed{\begin{aligned}\deleted{\unsure{Q((x_1+u_1)\wedge(x_2+u_2))}}&\\[-2pt]\deleted{\unsure{Q((x_1+u_1)*(x_2+u_2))}}\quad Q(u_1*u_2)&=b_{\unsure{1|2}}^2Q_1(\deleted{\unsure{x_1+}}u_1)+b_2^2Q_2(\deleted{\unsure{x_2+}}u_2)\\ &\quad+4\big[Q_1(u_1)Q_2(u_2)-c_1Q_2(u_2)-c_2Q_1(u_1)\big]\end{aligned}}}
\]
% manuscript line(s) 21
Faisant $u_1=0$, $u_2=0$, on trouve $u=0$ d'où
% manuscript line(s) 22–23
\[
\underset{\substack{\shortparallel\\ c\deleted{\unsure{x}}}}{Q(x_1*x_2)}=b_1^2c_2+b_2^2c_1+\underbrace{4(c_1c_2-c_1c_2-c_1c_2)}_{-4c_1c_2}
\]
% manuscript line(s) 24
\deleted{\unsure{$Q(x_1,x_2)$}} La formule déjà donnée.
% archive sheet number: 64
\scanpage{66}

% manuscript line(s) 1
Retenons donc
% manuscript line(s) 2–3
\[
\boxed{\underset{\substack{\shortparallel\\c}}{Q(x_1*x_2)}=\underset{\substack{\shortparallel\\c_1}}{Q_1(x_1)}\delta_2+\underset{\substack{\shortparallel\\c_2}}{Q_2(x_2)}\delta_1+4Q_1(x_1)Q_2(x_2)}
\]
% manuscript line(s) 4
\[
c_1\delta_2+c_2\delta_1+4c_1c_2
\]

% manuscript line(s) 5
\uline{Digression affine}.
% manuscript line(s) 6
Soit $(X,\mathcal O)$ \unsure{topos|espace} ou topos annelé, muni d'une section
% manuscript line(s) 7
\[
\chi\in\Gamma\mathcal O
\]
% manuscript line(s) 8
On s'intéresse à la catégorie des triples
% manuscript line(s) 9
\[
(\deleted{\unsure{M}}L,P,T),
\]
% manuscript line(s) 10–12
où \quad $L$ un $O$-Module \quad $P$ un $L$-torseur \quad $T$ une trivialisation de ${}_{\unsure{T|\chi}}P\overset{\mathrm{def}}=P\wedge_L(L,L\overset{\chi}{\longrightarrow}L)$,
% manuscript line(s) 13
ou encore la catégorie des \deleted{triples} \unsure{système|systèmes}
% manuscript line(s) 14
\[
(M,E,\deleted{\unsure{q}}i,\psi,T)
\]
% manuscript line(s) 15
où $M,E$ des $O$-modules,
% manuscript line(s) 16
\[
0\unsure{\mathbin{-}|\to}M\overset{i}{\longrightarrow}E\overset{\psi}{\longrightarrow}\mathcal O\longrightarrow0\tag*{($*$)}
\]
% manuscript line(s) 16-right
\margin{\uline{NB} On va identifier $M\subset E$}
% manuscript line(s) 17–18
une suite exacte i.e. $E$ \deleted{\unsure{est}} \inserted{ét.} donné comme ext. de $O$ par $M$, et enfin $T\in E$ tel que
% manuscript line(s) 19
\[
\psi(T)=\chi
\]

% manuscript line(s) 20
(\uline{NB} On a ${}_{T}P=\psi^{-1}(\unsure{\pi|T|\chi})\subset E$).

% manuscript line(s) 21–22
On veut définir un ``produit contracté'' $\wedge$ sur cette catégorie
% manuscript line(s) 23
\[
(L_1,P_1,T_1)\wedge(L_2,P_2,T_2)=(L_1\otimes L_2,P_1\wedge P_2,\underset{\substack{\shortparallel\,?\\T_1\wedge T_2}}{T_{\unsure{*|1}}})
\]
% manuscript line(s) 24–25
qui fasse de cette catégorie une \deleted{\unsure{gr}} catégorie monoïdale (contraintes d'associativité, d'unité, de comm.).
% archive sheet number: 65
